\documentclass[12pt]{article}
\usepackage{ctex}
\usepackage{sci-manuscript}
\addbibresource{references.bib}

\ManuscriptTitle{\LaTeX 模板标题}

\ManuscriptAuthors{%
  \AuthorORCID{Lintao Luo}{1,*}{0009-0001-5790-4866}
}

\ManuscriptAffiliations{%
  \Affiliation{1}{giantlinlinlin@gmail.com}%
}

\Keywords{关键词1; 关键词2; 关键词3}

% \Correspondence{*Corresponding author: author@example.edu}

\begin{document}

\MakeManuscriptTitle

\begin{manuscriptabstract}
这是一个方便人工智能编辑的\LaTeX 模板，可以直接用于撰写科学论文。该模板支持中文和英文混排。撰写时把本段替换为真正的摘要。
\end{manuscriptabstract}

\section{Results}
在这里撰写你的SCI结果部分。必要时引用参考文献，例如\parencite{aton2025}。


\section{Materials and Methods}

\subsection{Population structure}
Population structure was inferred using \software{ADMIXTURE} (v1.3.0)
for ancestral population numbers ranging from
$K = 2$ to $K = 10$.

\subsection{Coalescence time estimation}
The coalescence time of each lineage was estimated using a $\rho$ statistic-based
method \parencite{forster1996} and an ML method \parencite{behar2012} with the
Soares rate for complete mtDNA sequences (\SI{3624}{years} per substitution,
corrected for purifying selection) \parencite{soares2009}.

\printbibliography[title={References}]

\end{document}
